\chapter{REFERENCIAL TEÓRICO}
\label{cap:referencial}

\section{Geração Procedural de Conteúdo em Jogos Digitais}
\label{sec:pcg}

A geração procedural de conteúdo, conhecida pela sigla PCG
(\textit{Procedural Content Generation}), corresponde ao uso de
algoritmos e procedimentos computacionais para produzir
automaticamente, ou com algum grau de intervenção humana,
elementos que compõem um jogo digital. Esses elementos podem
incluir fases, mapas, cenários, personagens, itens, regras
e outros componentes necessários à experiência de jogo.
Hendrikx et al. (2013) apresentam uma classificação dos
diferentes tipos de conteúdo que podem ser gerados
proceduralmente, organizando-os em categorias como espaço,
sistemas, cenários e design. Complementarmente, a coletânea 
organizada por Shaker, Togelius e Nelson (2016) reúne discussões 
sobre métodos e aplicações da geração procedural no desenvolvimento 
de jogos digitais \cite{Hendrikx2013,Shaker2016,TogeliusSearch2016}.

A utilização da geração procedural está relacionada a
diferentes objetivos no desenvolvimento de jogos. Entre
eles, destacam-se a possibilidade de produzir grandes
quantidades de conteúdo, ampliar a variedade de experiências
oferecidas ao jogador, aumentar a reaproveitabilidade dos
elementos de um jogo e auxiliar o trabalho de designers.
Entretanto, a geração automática não significa,
necessariamente, que o conteúdo produzido seja adequado
ou interessante. Um algoritmo pode gerar uma grande
quantidade de fases, mas ainda assim produzir estruturas
repetitivas, desafios desbalanceados ou elementos
incompatíveis com as regras do jogo. Portanto, a qualidade
do conteúdo gerado constitui uma questão central no
desenvolvimento de sistemas de PCG
\cite{Hendrikx2013,Shaker2016,MalekiZhao2024}.

Os métodos de PCG podem ser diferenciados pela forma
como produzem o conteúdo. Nos métodos construtivos,
a produção ocorre diretamente por procedimentos
previamente definidos, sem que seja necessário avaliar 
repetidamente o resultado para gerar novas alternativas. 
Já os métodos baseados em busca produzem e avaliam diferentes 
soluções candidatas, utilizando critérios de avaliação para 
orientar a procura por resultados que atendam a determinados
objetivos. Sistemas também podem utilizar regras e
mecânicas para organizar os elementos do jogo. Essas
abordagens não são necessariamente excludentes e
podem ser combinadas de diferentes maneiras
\cite{ShakerConstructive2016,TogeliusSearch2016,
Togelius2011,NelsonRules2016}.

Além dessas abordagens, a PCG também pode utilizar
técnicas de aprendizado de máquina. Nesse caso, modelos
computacionais são treinados a partir de conteúdos
existentes para identificar padrões e produzir novos
elementos com características semelhantes às dos dados
utilizados no treinamento. Summerville et al. (2018)
apresentam uma revisão sobre a geração procedural por
aprendizado de máquina, discutindo diferentes modelos
e suas aplicações na criação de conteúdo funcional
para jogos, incluindo fases de jogos de plataforma.
Essa abordagem representa uma alternativa aos métodos
construtivos e de busca, mas depende de dados adequados
para o treinamento e apresenta desafios relacionados
à representação e à avaliação do conteúdo produzido
\cite{Summerville2018}.

Mais recentemente, Maleki e Zhao (2024) discutem a
evolução das técnicas de PCG, abordando métodos baseados
em busca, aprendizado de máquina, funções de ruído
e modelos de linguagem de grande escala. Os autores
também discutem possibilidades de combinação entre
diferentes técnicas. Essa diversidade demonstra que
a geração procedural não está limitada a um único
tipo de algoritmo, mas compreende um conjunto de
estratégias que podem ser selecionadas de acordo
com o tipo de conteúdo e os objetivos do sistema 
\cite{MalekiZhao2024}.

No contexto dos jogos de plataforma 2D, a geração procedural de fases 
apresenta desafios que vão além da produção de estruturas visualmente 
distintas. É necessário considerar se as fases geradas podem ser 
efetivamente completadas pelo jogador, respeitando as mecânicas e 
regras estabelecidas pelo jogo. Nesse sentido, aspectos como 
completabilidade, diversidade e adequação dos desafios constituem 
características importantes para a avaliação da qualidade do conteúdo 
gerado. Além disso, a análise do desempenho dos algoritmos de geração 
e da percepção dos jogadores permite considerar diferentes perspectivas 
na avaliação das fases produzidas.

\subsection{Avaliação da qualidade do conteúdo gerado}

A avaliação é uma etapa necessária para compreender se
um sistema de geração procedural atende aos objetivos
para os quais foi desenvolvido. A quantidade de conteúdo
produzido, isoladamente, não é suficiente para determinar
a qualidade de um gerador. É necessário considerar
características relacionadas à funcionalidade, à variedade
e à adequação do conteúdo ao contexto do jogo
\cite{ShakerSmithYannakakis2016,Khalifa2025}.

Khalifa et al. (2025) apresentam o
\textit{Procedural Content Generation Benchmark}, uma
estrutura de avaliação que reúne diferentes problemas
de geração de conteúdo para jogos. O trabalho considera
três dimensões principais: qualidade, diversidade e
controlabilidade. A qualidade está relacionada ao
atendimento dos critérios funcionais definidos para
o problema; a diversidade considera a variedade dos
conteúdos produzidos; e a controlabilidade avalia se
o gerador consegue respeitar parâmetros de controle
estabelecidos. Os autores demonstram que diferentes
algoritmos podem apresentar desempenhos distintos
conforme o problema e os objetivos de geração
\cite{Khalifa2025}.

Outra abordagem de avaliação é a análise do espaço
de possibilidades que um gerador consegue produzir,
frequentemente denominada \textit{Expressive Range
Analysis} (ERA). Essa análise busca compreender quais
tipos de conteúdo são produzidos, quais características
variam e quais limitações podem estar presentes no
algoritmo. Smith e Whitehead (2010), por exemplo,
aplicaram essa perspectiva ao gerador Launchpad,
desenvolvido para produzir fases de jogos de plataforma
2D. Os autores utilizaram características como
linearidade e leniência dos níveis para analisar
a variedade dos resultados e identificar tendências
e limitações do gerador \cite{SmithWhitehead2010}.

A avaliação também pode envolver a participação de
jogadores. Nesse caso, os participantes interagem com
o conteúdo gerado e fornecem informações sobre aspectos
que nem sempre podem ser representados adequadamente
por métricas computacionais, como diversão, satisfação
e dificuldade percebida. Shaker, Smith e Yannakakis
(2016) discutem a utilização de visualizações do espaço
expressivo e questionários para compreender diferentes
aspectos da qualidade dos geradores. De maneira
complementar, Summerville et al. (2018) investigam
a relação entre métricas computacionais de fases
de plataforma e avaliações humanas. Essas abordagens
podem ser complementares: as métricas automatizadas
permitem avaliar grandes quantidades de conteúdo,
enquanto a avaliação humana contribui para compreender
como os jogadores percebem e experimentam esse conteúdo
\cite{ShakerSmithYannakakis2016,Summerville2018}.

Além das características de qualidade e diversidade,
a avaliação de sistemas de geração pode considerar
os recursos computacionais necessários para produzir
e avaliar os conteúdos. O tempo de geração e o custo
das avaliações realizadas durante a busca, por exemplo,
permitem analisar o desempenho computacional dos
geradores. Withington, Cook e Tokarchuk (2024) incluem
o uso de recursos computacionais entre os aspectos
observados nas práticas de avaliação de sistemas
de geração procedural \cite{Withington2024}.

No contexto de fases de plataforma, é importante
separar a validade estrutural da possibilidade de
conclusão. Neste trabalho, a completabilidade refere-se,
conceitualmente, à existência de uma sequência de ações
que permita ao personagem alcançar o objetivo final
a partir do ponto inicial, respeitando o modelo de
movimentação e colisão do jogo. Essa propriedade não
representa, isoladamente, a qualidade da experiência
oferecida ao jogador.

Dessa maneira, a avaliação de um gerador procedural
deve considerar os objetivos específicos da aplicação.
Para a geração de fases de plataforma 2D, é importante
verificar se as estruturas respeitam as restrições
do jogo, se podem ser completadas
segundo o modelo de movimentação considerado e se
os conteúdos produzidos apresentam diferenças
relevantes entre si. Além disso, a avaliação com jogadores
permite investigar como os conteúdos são percebidos
durante a interação. Essas dimensões são complementares
e não devem ser tratadas como medidas equivalentes.

\section{Geração Procedural de Fases em Jogos de Plataforma 2D}
\label{sec:geracao-fases}

A geração procedural de fases consiste na utilização de
métodos computacionais para construir ambientes e desafios
que serão percorridos pelo jogador durante uma partida.
Em jogos de plataforma 2D, esses ambientes podem ser
compostos por terrenos, plataformas, espaços vazios,
obstáculos e pontos de início e término. A disposição
desses elementos influencia as possibilidades de
movimentação do personagem e os desafios apresentados
ao jogador.

A geração de fases nesse gênero apresenta desafios
relacionados à interação entre a estrutura do ambiente
e as mecânicas de movimentação. A distância horizontal
entre plataformas, as diferenças de altura, a presença
de obstáculos e o espaço disponível para saltos podem
determinar se uma região é alcançável. Dessa forma,
uma fase que apresenta uma organização geométrica
aparentemente válida não é necessariamente completável.

Essa distinção é importante porque a geração de uma
estrutura e a verificação de sua jogabilidade representam
problemas diferentes. A estrutura pode respeitar limites
de dimensão e posicionamento, mas ainda apresentar
uma sequência de plataformas que não pode ser alcançada
considerando a velocidade, a gravidade e as características
do salto do personagem. Portanto, a avaliação de fases
de plataforma precisa considerar as regras que determinam
a movimentação e a interação com o ambiente.

Summerville et al. (2017) apresentam métricas para
descrever características de fases de plataforma e
investigam sua relação com avaliações humanas. O estudo
considera aspectos como dificuldade percebida, estética
visual e diversão, demonstrando a importância de
relacionar características estruturais das fases com
a experiência dos jogadores. Essa abordagem contribui
para compreender que a avaliação de uma fase não deve
se limitar à sua representação geométrica
\cite{Summerville2017}.

A construção de fases também pode utilizar restrições
e estruturas previamente definidas. Shaker et al. (2016)
apresentam métodos construtivos para geração de níveis,
nos quais procedimentos e condições determinam como
os elementos são organizados. Dormans (2010), por sua
vez, discute o uso de gramáticas generativas para
construir missões e espaços em jogos de ação e aventura.
Embora o domínio investigado pelo autor não seja
exatamente o mesmo deste trabalho, sua discussão
contribui para compreender como regras formais podem
organizar estruturas de conteúdo
\cite{ShakerConstructive2016,Dormans2010}.
\subsection{Estratégias de geração de fases}

A geração procedural de fases pode utilizar diferentes
estratégias para determinar como os elementos de um
cenário são selecionados e organizados. Entre as
abordagens discutidas na literatura estão os métodos
construtivos e os métodos baseados em busca. Nos métodos
construtivos, o conteúdo é produzido diretamente por
procedimentos de construção; nos métodos baseados em
busca, soluções candidatas são geradas e avaliadas
para orientar a procura por configurações que atendam
a determinados objetivos. Essa distinção permite compreender 
diferentes formas de organização dos processos de geração, embora 
as abordagens não sejam necessariamente excludentes 
\cite{ShakerConstructive2016,TogeliusSearch2016,Togelius2011}.

\subsection{Geração aleatória}

A geração aleatória utiliza sorteios para determinar
características ou componentes do conteúdo produzido.
Em sistemas de geração procedural, esses sorteios
podem selecionar elementos, valores ou configurações
dentro de um conjunto de possibilidades previamente
definido. Maleki e Zhao (2024) apresentam os geradores
de números pseudoaleatórios como uma técnica utilizada
em PCG e explicam que sua operação envolve uma semente
(\textit{seed}), um algoritmo e uma distribuição dos
valores produzidos. Embora as sequências geradas sejam
pseudoaleatórias, a utilização da mesma semente e do
mesmo algoritmo permite reproduzir os resultados
\cite{MalekiZhao2024}.

No domínio considerado nesta pesquisa, a aleatoriedade
pode ser empregada para selecionar características
como a largura dos segmentos, a altura das plataformas,
a distância entre elas e a presença ou o tipo de
obstáculos. Os valores podem ser sorteados dentro de
intervalos definidos para o domínio do jogo. A utilização
desses limites restringe as possibilidades de geração,
mas não elimina o caráter aleatório da escolha: dentro
do conjunto permitido, os valores continuam sendo
determinados pelo mecanismo de sorteio. Métodos
construtivos podem incorporar processos estocásticos
para selecionar e organizar componentes durante a
produção de níveis \cite{ShakerConstructive2016}.

Uma característica desse tipo de abordagem é que
ela permite produzir diferentes configurações a partir
de um mesmo conjunto de parâmetros. Além disso, quando
o gerador utiliza uma semente registrada, uma execução
pode ser reproduzida, o que é relevante para testar
fases específicas e comparar resultados experimentais
\cite{MalekiZhao2024}. Entretanto, a aleatoriedade,
por si só, não orienta a geração para configurações
consideradas melhores segundo um objetivo de qualidade.
A produção de conteúdos adequados depende dos parâmetros,
das possibilidades disponibilizadas e das verificações
aplicadas durante ou após a construção.

Dessa forma, a geração aleatória permite produzir diferentes configurações 
a partir de sorteios realizados sobre um conjunto de possibilidades 
previamente definido. Embora esse mecanismo favoreça a variedade das 
configurações produzidas, a aleatoriedade, isoladamente, não garante que 
os conteúdos atendam a critérios específicos de qualidade ou 
funcionalidade. Por esse motivo, sua utilização pode ser acompanhada por 
procedimentos de validação e avaliação.

\subsection{Geração baseada em regras}

A geração baseada em regras utiliza condições explícitas
para orientar a construção e a organização dos elementos
de um conteúdo. Essas condições podem estabelecer quais
componentes podem ser utilizados, em que situações
podem aparecer e como devem se relacionar com os
componentes já produzidos. Nos métodos construtivos,
procedimentos definidos previamente podem determinar
a sequência de construção e a seleção dos elementos
que formarão um nível \cite{ShakerConstructive2016}.

Como exemplo de aplicação ao domínio desta pesquisa,
uma regra pode orientar a escolha do próximo segmento considerando
a altura da plataforma anterior, a distância entre
plataformas ou a presença de determinado obstáculo.
Também podem ser estabelecidas condições para evitar
combinações estruturais indesejadas. Nesse tipo de
abordagem, as regras participam das decisões de
construção: elas não apenas verificam se o resultado
final respeita limites, mas influenciam quais elementos
serão selecionados e como serão organizados.

É importante distinguir o uso de regras na construção
de um conteúdo da aplicação de restrições gerais de
validade. Um gerador não se caracteriza como baseado
em regras apenas por limitar dimensões ou rejeitar
configurações inválidas. Para que as regras orientem
a geração, suas condições precisam participar do
processo de decisão, por exemplo, determinando quais
segmentos podem suceder um segmento anterior. Essa
compreensão é compatível com a discussão dos métodos
construtivos de geração de níveis apresentada por
Shaker et al. (2016) \cite{ShakerConstructive2016}.

Uma abordagem relacionada é a utilização de gramáticas
generativas. Embora também empreguem regras explícitas,
gramáticas possuem uma estrutura formal de produção,
na qual símbolos e regras de substituição permitem
derivar configurações. Togelius, Shaker e Dormans
(2016) discutem gramáticas e sistemas L e apresentam
aplicações na geração de vegetação e de níveis para
jogos. Portanto, gramáticas podem ser utilizadas como
uma forma estruturada de geração baseada em regras,
mas nem todo gerador que utiliza condições explícitas
precisa implementar uma gramática formal
\cite{TogeliusGrammars2016}.

O uso de regras pode contribuir para organizar os
componentes e restringir combinações incompatíveis
com as condições definidas pelo desenvolvedor.
Entretanto, a aplicação dessas condições não garante,
por si só, que a fase possa ser completada. Uma
configuração pode respeitar todas as regras estruturais
e ainda apresentar uma sequência de plataformas
inacessível de acordo com a velocidade, a gravidade
e as capacidades de salto do personagem. Assim,
a validade estrutural e a completabilidade precisam
ser verificadas por procedimentos distintos.

A geração baseada em regras permite incorporar condições explícitas ao 
processo de construção, orientando a seleção e a organização dos 
componentes do conteúdo. Essa abordagem oferece maior controle sobre 
determinadas características estruturais, mas a utilização de regras não 
garante, por si só, que o conteúdo produzido seja funcional ou completável. 
A adequação das condições estabelecidas depende das características do 
domínio e dos objetivos do sistema de geração.

\section{Algoritmos Genéticos Aplicados à Geração de Fases}
\label{sec:algoritmos-geneticos}

Os algoritmos genéticos são métodos de busca e otimização
pertencentes à classe dos algoritmos evolutivos. Inspirados
em mecanismos da evolução biológica, utilizam processos
como seleção, cruzamento e mutação para explorar possíveis
soluções de um problema. Seu funcionamento é baseado
em uma população de indivíduos, na qual cada indivíduo
representa uma solução candidata. A qualidade dessas
soluções é estimada por uma função de aptidão, também
denominada \textit{fitness}, que permite orientar o processo
evolutivo em direção aos objetivos estabelecidos
\cite{Goldberg1989,Mitchell1996,EibenSmith2003}.

De maneira geral, o processo começa com a criação de uma
população inicial de soluções candidatas. Em seguida,
os indivíduos são avaliados de acordo com uma função
de aptidão, e alguns deles são selecionados para participar
da reprodução. Operadores genéticos são aplicados aos
indivíduos selecionados, produzindo descendentes que
podem compor uma nova população. Esse processo ocorre
ao longo de sucessivas gerações e é encerrado quando
um critério de parada é atingido, como um número máximo
de gerações ou a obtenção de uma solução que satisfaça
determinada condição \cite{Goldberg1989,EibenSmith2003}.

A aplicação desses algoritmos depende de diferentes
elementos, entre os quais se destacam a representação
das soluções, a função de aptidão e os operadores
genéticos. Esses componentes estão relacionados,
pois a forma como uma solução é representada influencia
as modificações que podem ser realizadas, enquanto
a função de aptidão determina quais características
são favorecidas durante a evolução da população
\cite{Mitchell1996,EibenSmith2003}.

\subsection{Representação dos indivíduos e função de aptidão}

A representação dos indivíduos, também denominada
codificação, define como as soluções candidatas são
estruturadas e armazenadas para que possam ser
manipuladas pelo algoritmo. Diferentes formas de
representação podem ser utilizadas, como sequências,
matrizes, árvores e grafos. A escolha depende das
características do problema e da estrutura das soluções
que se deseja produzir \cite{Goldberg1989,Togelius2011,EibenSmith2003}.

Em problemas de geração procedural de fases, a
representação precisa descrever os elementos que
compõem o cenário e suas relações. Uma representação
em matriz, por exemplo, permite associar elementos
a posições específicas de uma estrutura bidimensional.
Já uma representação sequencial organiza os componentes
em uma ordem, permitindo representar conteúdos formados
por partes consecutivas. A representação escolhida
também influencia a aplicação dos operadores genéticos,
pois determina quais modificações podem ser realizadas
e como os indivíduos podem ser combinados.

A função de aptidão é responsável por atribuir uma
avaliação às soluções candidatas. Seu propósito é
indicar o quanto cada indivíduo atende aos critérios
considerados relevantes para o problema, orientando
a seleção durante o processo evolutivo. Em sistemas
de geração de fases, esses critérios podem estar
relacionados a características como completabilidade,
dificuldade, organização espacial e outras propriedades
do conteúdo produzido \cite{Goldberg1989,Togelius2011}.

A definição da função de aptidão exerce influência
direta sobre o comportamento do algoritmo, pois
características que recebem maior valorização tendem
a ser favorecidas durante a busca. Consequentemente,
uma função que considera apenas determinados aspectos
pode levar à produção de soluções que apresentam
bons resultados nesses critérios, mas não necessariamente
atendem a outras expectativas relacionadas à qualidade
do conteúdo. Por esse motivo, a função de aptidão
deve ser compreendida como um mecanismo de orientação
da busca, e não como uma representação completa
de todas as dimensões da qualidade de uma solução
\cite{Mitchell1996,EibenSmith2003}.

\subsection{Operadores genéticos}

Os operadores genéticos são responsáveis por
selecionar indivíduos e produzir novas soluções
a partir da população. Entre os operadores
tradicionais estão a seleção, o cruzamento e
a mutação. A combinação desses mecanismos permite explorar
diferentes regiões do espaço de soluções e transmitir
características dos indivíduos ao longo das gerações
\cite{Goldberg1989,Mitchell1996,EibenSmith2003}.

A seleção determina quais indivíduos terão a
oportunidade de participar do processo de reprodução.
Um dos métodos utilizados é a seleção por torneio,
na qual um conjunto de indivíduos é escolhido
e aquele que apresenta a melhor aptidão dentro
desse grupo é selecionado. O procedimento pode
ser repetido até que seja obtida a quantidade
necessária de indivíduos. O tamanho do grupo
participante influencia a pressão seletiva:
torneios maiores tendem a favorecer indivíduos
com maior aptidão, enquanto torneios menores
podem permitir maior participação de soluções
com avaliações diferentes \cite{EibenSmith2003}.

O cruzamento, também denominado recombinação,
combina partes de dois indivíduos para produzir
novas soluções candidatas. No cruzamento em um
ponto, uma posição da representação é selecionada,
e os componentes dos indivíduos são divididos
a partir desse ponto para formar os descendentes.
Essa operação permite combinar características
presentes nas soluções utilizadas como progenitoras.
A aplicação do cruzamento depende da estrutura
da representação, pois a combinação precisa
preservar uma interpretação válida dos indivíduos
\cite{Goldberg1989,Mitchell1996}.

A mutação introduz alterações nos indivíduos,
permitindo explorar configurações que não estejam
presentes na população atual. Sua aplicação pode
envolver a alteração de valores, a substituição
de elementos ou modificações na estrutura da
solução. A mutação contribui para manter a
variabilidade da população e ampliar a exploração
do espaço de soluções, embora sua intensidade
precise ser considerada em relação ao comportamento
geral do algoritmo \cite{EibenSmith2003}.

Outra estratégia utilizada em algoritmos genéticos
é o elitismo, que consiste em preservar determinados
indivíduos com melhor aptidão para que sejam
transferidos diretamente à geração seguinte.
Esse mecanismo reduz a possibilidade de perda
de soluções bem avaliadas em decorrência das
operações de cruzamento e mutação. Entretanto,
a preservação excessiva dos melhores indivíduos
pode reduzir a diversidade da população e favorecer
a concentração da busca em determinadas regiões
do espaço de soluções \cite{Goldberg1989,EibenSmith2003}.

A aplicação dos operadores também pode produzir
soluções que não respeitam as restrições do problema
ou da representação utilizada. Em determinadas
abordagens, são empregados mecanismos de validação
ou reparação para identificar e corrigir essas
inconsistências. A necessidade desses mecanismos
está relacionada à forma de codificação e às
características das operações realizadas. No caso
da geração de fases, a validade estrutural de
uma solução não implica necessariamente que
o conteúdo produzido seja funcional ou completável.

\subsection{Aplicação de algoritmos genéticos à geração de fases}

A utilização de algoritmos genéticos na geração
procedural de conteúdo permite explorar diferentes
configurações de um artefato por meio da evolução
de soluções candidatas. Nesse contexto, as fases
podem ser representadas como indivíduos de uma
população, enquanto a função de aptidão avalia
características relacionadas aos objetivos da
geração. A representação e os critérios de avaliação
são, portanto, aspectos fundamentais para determinar
quais soluções podem ser exploradas e quais
características serão favorecidas durante o
processo evolutivo \cite{Togelius2011}.

Togelius et al. (2011) discutem a aplicação de
métodos evolutivos à geração procedural baseada
em busca, destacando a relação entre representação,
avaliação e exploração do espaço de soluções.
Essa abordagem permite orientar a produção de
conteúdo de acordo com critérios definidos,
mas sua efetividade depende da adequação desses
critérios ao domínio em que o gerador será utilizado
\cite{Togelius2011}.

Sabraw e Churchill (2022) apresentam um sistema
baseado em algoritmo genético para produzir fases
de um jogo de plataforma 2D. O estudo é retomado
na seção de trabalhos relacionados, na qual são
discutidas sua representação, sua função de
avaliação e as diferenças em relação ao escopo
deste trabalho \cite{SabrawChurchill2022}.

Apesar das possibilidades oferecidas por essa
abordagem, a utilização de algoritmos genéticos
apresenta desafios. O processo evolutivo depende
da avaliação repetida de soluções candidatas,
o que pode aumentar o custo computacional,
especialmente quando a avaliação exige a execução
de simulações. Além disso, a escolha da representação,
dos operadores e dos parâmetros influencia a
exploração do espaço de soluções e a diversidade
dos indivíduos ao longo das gerações
\cite{Togelius2011,EibenSmith2003}.

Assim, os algoritmos genéticos constituem uma
abordagem de busca que permite explorar e modificar
soluções candidatas de maneira evolutiva. Sua
aplicação à geração de fases depende da relação
entre a representação do conteúdo, os critérios
de avaliação e os operadores utilizados. Esses
elementos determinam características importantes
do funcionamento do gerador e ajudam a compreender
tanto suas possibilidades quanto suas limitações.

\section{Trabalhos Relacionados}
\label{sec:trabalhos-relacionados}

Esta seção apresenta trabalhos que investigam a geração
procedural de fases, a avaliação de conteúdos de jogos
de plataforma e a comparação entre sistemas de geração.
A análise considera os objetivos de cada estudo, as
técnicas utilizadas, os principais resultados e as
diferenças em relação à proposta desta pesquisa.

\subsection{Geração de fases e análise de características estruturais}

Sabraw e Churchill (2022), no trabalho \textit{Generating
2D Platforming Auto-Play Levels with Genetic Algorithms},
investigam a utilização de algoritmos genéticos para gerar 
fases de um jogo de plataforma 2D em um motor desenvolvido 
pelos próprios autores. O sistema utiliza diferentes elementos 
de cenário, como blocos convencionais, esteiras, superfícies 
de gelo e blocos que provocam impulsos no personagem. 
As fases são representadas em uma estrutura de grade, que 
permite organizar os elementos e aplicar operações durante 
o processo evolutivo \cite{SabrawChurchill2022}.

Os autores investigam diferentes configurações de
geração, incluindo a aleatorização de elementos,
alterações de parâmetros e restrições no espaço
disponível para a construção das fases. A avaliação
considera o deslocamento do personagem para a direita,
e os experimentos também investigam uma função de
aptidão denominada \textit{least-blocks}, proposta
para reduzir a quantidade de blocos considerados excessivos.
Os resultados indicam que essa alteração
pode modificar a densidade das fases, produzindo tanto
configurações menos carregadas quanto cenários
excessivamente escassos em determinadas condições
\cite{SabrawChurchill2022}.

A principal contribuição desse estudo para esta
pesquisa está na demonstração de que algoritmos
genéticos podem ser aplicados à geração de fases
de plataforma 2D, utilizando uma representação
computacional e critérios de avaliação automatizada.
Entretanto, as fases produzidas por Sabraw
e Churchill (2022) são do tipo \textit{auto-play},
nas quais o personagem não é controlado diretamente
pelo jogador. No presente trabalho, pretende-se gerar
fases para um jogo em que o participante controla
o personagem, o que exige considerar as ações
de movimentação e salto durante a avaliação de
completabilidade \cite{SabrawChurchill2022}.

Schaa e Barriga (2024), no artigo \textit{Evaluating the
Expressive Range of Super Mario Bros Level Generators}, 
apresentam uma avaliação do espaço expressivo de três 
geradores automáticos de fases de \textit{Super Mario Bros.}. 
Os sistemas investigados utilizam diferentes abordagens:
algoritmo genético com múltiplas populações, uma
rede adversarial generativa convolucional profunda
(DCGAN) combinada com uma estratégia de adaptação
de matriz de covariância (CMA-ES) e cadeias de
Markov. Os autores também consideram as fases
originais utilizadas como referência para a
comparação \cite{SchaaBarriga2024}.

A avaliação utiliza a técnica de
\textit{Expressive Range Analysis} (ERA), com
nove métricas relacionadas à estrutura e à
dificuldade das fases. Foram geradas 200 fases
para cada gerador, permitindo analisar a
distribuição das características dos conteúdos
produzidos. Os resultados indicam que os algoritmos 
genéticos e as cadeias de Markov apresentam espaços 
expressivos mais amplos que o gerador baseado em GAN. 
Os autores também observam que o algoritmo genético 
apresenta diversidade, apesar da convergência antecipada, 
enquanto o gerador baseado em GAN produz, em algumas 
situações, conteúdos pouco coerentes e com diversidade 
limitada \cite{SchaaBarriga2024}.

Esse estudo contribui para a presente pesquisa por
mostrar como métricas estruturais podem ser utilizadas
para comparar a diversidade e as características de
fases de plataforma. Entretanto, seu foco está na análise
do espaço expressivo de geradores existentes para 
\textit{Super Mario Bros.}, enquanto este presente 
trabalho pretende implementar e comparar três mecanismos 
de geração em um mesmo protótipo, incluindo a avaliação 
de completabilidade e a participação de jogadores.

\subsection{Métricas computacionais e avaliação da experiência dos jogadores}

Summerville et al. (2017), no trabalho \textit{Understanding
Mario: An Evaluation of Design Metrics for Platformers},
investigam métricas computacionais para avaliar
características de fases de jogos de plataforma e sua
relação com percepções humanas. O estudo utiliza fases
do jogo \textit{Infinite Mario Bros.} e analisa se as
métricas propostas conseguem representar aspectos como
dificuldade percebida, estética visual e diversão.
Para isso, os autores comparam os valores calculados
pelas métricas com avaliações fornecidas por participantes
\cite{Summerville2017}.

Os resultados demonstram que algumas métricas apresentam
correlações com as avaliações humanas. A quantidade
de inimigos, por exemplo, alcançou correlação de 0,72
com a dificuldade percebida e de 0,42 com a diversão.
Para a estética visual, a maior correlação reportada
foi de 0,23. Esses resultados indicam que determinadas
características computacionais podem contribuir para
explicar aspectos da experiência, mas não representam
integralmente as percepções dos jogadores
\cite{Summerville2017}.

A relação desse trabalho com o presente projeto de 
pesquisa está na utilização de características mensuráveis 
das fases para apoiar sua avaliação e na consideração 
da percepção humana. Entretanto, Summerville et al. (2017) 
concentram-se na análise de métricas e em sua relação com
avaliações humanas, enquanto esta pesquisa pretende comparar 
três estratégias de geração procedural em um protótipo próprio.
As métricas serão utilizadas para analisar os conteúdos
produzidos, em conjunto com a verificação de completabilidade 
e a avaliação dos participantes.

\subsection{Avaliação sistemática de geradores e benchmarks de PCG}

Withington, Cook e Tokarchuk (2024), no trabalho
\textit{On the Evaluation of Procedural Level Generation
Systems}, investigam as práticas de avaliação utilizadas em
sistemas de geração procedural de fases. Os autores
desenvolvem uma taxonomia para organizar essas práticas
e analisam publicações acadêmicas, buscando identificar 
como os pesquisadores coletam dados, caracterizam os
conteúdos produzidos e comparam seus sistemas 
\cite{Withington2024}.

A taxonomia organiza a avaliação em torno de quatro
questões: o método utilizado para coletar dados sobre
o sistema, as métricas e características calculadas,
o ponto de comparação adotado e o domínio de jogo
em que o sistema é avaliado. A pesquisa parte de um
conjunto de 137 publicações e seleciona 86 trabalhos
que apresentam sistemas novos de geração de fases.
A análise evidencia a variedade de práticas utilizadas
e fundamenta recomendações relacionadas à transparência
metodológica, à diversidade de estruturas de pesquisa
e ao reaproveitamento de códigos e metodologias
\cite{Withington2024}.

A contribuição desse trabalho para o presente estudo está
na fundamentação da necessidade de definir critérios de 
avaliação explícitos e adequados ao problema investigado. 
A proposta considera essa preocupação ao estabelecer três 
níveis de avaliação: validação estrutural, avaliação automatizada 
e avaliação com participantes, além de registrar as condições 
utilizadas na comparação.

Khalifa et al. (2025), no trabalho \textit{The Procedural
Content Generation Benchmark: An Open-source Testbed
for Generative Challenges in Games}, apresentam uma 
estrutura de código aberto desenvolvida para avaliar 
algoritmos de geração procedural em diferentes problemas 
relacionados a jogos. O benchmark reúne 12 problemas, cada 
um com múltiplas variantes, que abrangem desde a geração de 
fases até a criação de conjuntos de regras para jogos simples. 
Cada problema possui representação, parâmetros de controle
e critérios próprios de avaliação \cite{Khalifa2025}.

Os autores propõem avaliar os conteúdos produzidos considerando 
três dimensões: qualidade, diversidade e controlabilidade. O benchmark 
também é utilizado para comparar três algoritmos de referência: 
um gerador aleatório, uma estratégia evolutiva e um algoritmo genético. 
Os resultados demonstram que o desempenho dos algoritmos varia conforme 
o problema e os objetivos de avaliação estabelecidos. Os autores também 
observam que a geração de fases maiores, como as de \textit{Super Mario Bros.}
e \textit{Lode Runner}, representa um desafio para os algoritmos avaliados
\cite{Khalifa2025}.

Esse trabalho contribui para a presente pesquisa por apresentar uma 
referência metodológica para a comparação sistemática de diferentes
algoritmos de geração. Entretanto, seu escopo abrange diferentes problemas 
de conteúdo para jogos, enquanto este TCC concentra-se em fases de 
plataforma 2D e acrescenta a avaliação de completabilidade segundo o modelo 
de movimentação do personagem e a investigação da experiência dos 
participantes.

\subsection{Síntese dos trabalhos relacionados}

Os trabalhos analisados demonstram que a geração procedural de fases 
de plataforma pode ser investigada a partir de diferentes perspectivas. 
Sabraw e Churchill (2022) exploram o uso de algoritmos genéticos em
fases de plataforma \textit{auto-play}; Schaa e Barriga (2024) analisam o
espaço expressivo de diferentes geradores; Summerville et al. (2017) 
investigam métricas computacionais e sua relação com avaliações
humanas; Withington, Cook e Tokarchuk (2024) discutem as práticas de 
avaliação utilizadas na literatura; e Khalifa et al. (2025) apresentam
um benchmark para comparar algoritmos em diferentes problemas de geração 
\cite{SabrawChurchill2022,SchaaBarriga2024,Summerville2017,Withington2024,Khalifa2025}.

Em conjunto, esses estudos evidenciam a importância de avaliar sistemas 
de geração por mais de uma dimensão. A produção de fases estruturalmente 
válidas não garante que elas sejam completáveis, assim como a completabilidade
não representa, isoladamente, a qualidade da experiência oferecida ao jogador. 
Também é necessário considerar a diversidade dos conteúdos e o custo 
computacional envolvido na geração. Os trabalhos mostram ainda que métricas 
automatizadas e avaliações humanas oferecem informações complementares.

Diante desse contexto, esta pesquisa propõe comparar três mecanismos 
operacionais de geração procedural: aleatório, baseado em regras e por 
algoritmos genéticos, aplicados a um mesmo protótipo de jogo de plataforma 2D.
A comparação considera a completabilidade, a diversidade estrutural, o 
desempenho computacional dos geradores e a experiência dos jogadores. 
Busca-se analisar as diferenças entre essas estratégias em um domínio
experimental compartilhado, sem pressupor previamente a superioridade de 
qualquer uma delas.
